\section*{The Propers: Detailed Commentary}

\subsection*{The Sunday's lexical structure\enspace\properrefs{Int., Ep., All., Gosp., Sec.}}

The Missal numbers the appointed psalms by the Vulgate/Septuagint: 46, 33, and 30 (modern 47, 34, 31). Four lexical structures in the appointed texts carry much of the unity read across the formulary:

\begin{fourstable}{Words}{Immediate contrast}{Theological work}{Distinction}
\latin{voce} / \latin{manibus} & Voice and hands acclaim the King in Psalm 46. & Augustine directly interprets them as confession and good works in concord. & He comments on the psalm, not on this later Mass formulary.\\
\latin{exhibuistis} / \latin{exhibete} & The same members once presented to impurity are now to be presented to justice. & Conversion redirects embodied agency; freedom from one master becomes service to God. & Paul's slavery analogy is explicitly ``human'' language adapted to the hearers' weakness.\\
\latin{fructum} / \latin{fructus} & Romans names sanctification as fruit; Matthew names deeds as the fruit by which a tree is known. & The repeated noun makes holy conduct both a grace-produced reality and publicly relevant evidence. & A visible act does not disclose every intention or by itself prove doctrine true.\\
\latin{singuli} / \latin{cunctis} & Individuals offer; the Secret asks profit for all. & Personal oblation is ordered to ecclesial salvation within the one sacrifice. & The prayer does not dissolve persons into a collective or make every private intention universally efficacious by itself.\\
\end{fourstable}

\sectionguard
\subsection*{Psalm 46: a concordance of hands and voice\enspace\properrefs{Int., All.}}

Psalm 46 is an enthronement psalm: the nations are summoned to acclaim because God is the great King over all the earth. The Mass appoints its first command at both Introit and Alleluia and adds the reason at the entrance, \latin{Quoniam Dominus excelsus, terribilis: Rex magnus super omnem terram}. The repeated acclamation is therefore framed by sovereignty, not by applause as self-expression.

Augustine's \emph{Exposition on Psalm 46}, 2--3, supplies an unusually exact inherited bridge. He asks why the nations clap with hands and answers: because of good works. Mouth without active hands is defective; hands without confessing tongue are also defective. The two must agree. He then interprets the King over all the earth christologically: the one mocked and crucified has purchased the world by his blood and ascended as universal King.

\begin{threestable}{Liturgical occurrence}{What the text establishes}{What Augustine adds}
Introit & The whole assembly enters by a command addressed to all nations; hands and jubilant voice act together under the great King. & The bodily action signifies good works and the voice confession; Christ's blood grounds the worldwide people.\\
Alleluia & The same command is repeated immediately before Matthew's warning about false prophetic appearance and speech. & The inherited demand for agreement supplies a positive standard for the Gospel's negative warning.\\
Gospel response & Matthew contrasts saying ``Lord, Lord'' with doing the Father's will. & Augustine's psalm exposition and his Matthew exposition independently insist that a merely vocal confession is incomplete.\\
\end{threestable}

Augustine expounds the psalm and Matthew separately and directly; the later Mass, not he, sets the two texts side by side.

\sectionguard
\subsection*{Matthew 7: appearance, fruit, doctrine, and judgment\enspace\properrefs{Gosp.}}

The immediate context is essential. Jesus has just contrasted the narrow way ending in life with the broad way ending in destruction (Matt. 7:13--14). After the appointed verses he warns that prophecy, exorcism, and wonders done in his name do not replace righteousness (7:22--23), and he closes the sermon with two hearers: one hears and does, building on rock; the other hears and does not, building a ruin (7:24--27). The liturgical pericope stops at verse 21, but the complete paragraph shows why ``fruit'' cannot mean charisma, visibility, or religious vocabulary alone.

\begin{fourstable}{Direct witness}{Who is foregrounded?}{What counts as fruit?}{His caution}
St. John Chrysostom, \emph{Homilies 23--24 on Matthew} & In Homily 23 he chiefly distinguishes persons of corrupt life who mask themselves with virtue from open enemies; he does not primarily identify them with heretics. & The difficult, sustained practice of virtue exposes the show; in Homily 24, doing the Father's will outranks bare confession and even charismatic gifts. & An evil person may become good, but cannot bear good fruit while persisting in evil. Christ commands watchfulness, not private punishment.\\
St. Augustine, \emph{On the Sermon on the Mount} II.24--25 & He includes heretics and all deceptive seekers of worldly advantage under Christian appearance. & Manifest works of the flesh and fruit of the Spirit matter; trials involving hoped-for gain or feared loss disclose governing love. & Ascetic acts can be costume; hearts remain hidden. A bad teacher's true doctrine is the vine's grape gathered through thorns.\\
St. Jerome, \emph{Commentary on Matthew} I at 7:15--22 & He emphasizes heretics whose fasting, chastity, and restraint conceal poisonous teaching, then turns to orthodox speakers whose lives destroy their teaching. & Doctrine and conduct reciprocally confirm or contradict a claimant. & Moses, Jethro, Judas, and Paul show that ``tree'' is not an immutable human nature; repentance and fall remain possible.\\
\end{fourstable}

The differences are real. Chrysostom says that he thinks corrupt hypocrites rather than heretics are foremost; Augustine explicitly includes false claimants to Christian wisdom; Jerome gives heresy special emphasis and then equally insists on moral integrity. The three agree on conversion and on the insufficiency of appearance, but not on whom the warning chiefly targets.

St.~Anthony of Padua's \emph{Sermon for the Eighth Sunday after Pentecost} brings the same Gospel into a medieval preacher's pastoral examination. He directly coordinates wolves' concealment, the disclosure of fruit, and the difference between saying and doing, then turns the warning upon prelates and hearers rather than using it only against an external enemy. Anthony does not settle the Fathers' different identifications of the false prophets; his contribution is to make the text's test reflexive and ecclesial without claiming access to hidden hearts.

\begin{studybox}{Three distinct judgments}
\textbf{Moral fruit:} persistent, manifest action bears on credibility and may require avoidance. \textbf{Doctrinal norm:} teaching is measured by apostolic revelation, not generated by the teacher's attractiveness or vice. \textbf{Final judgment:} the fire and the kingdom belong to God; limited human judgment of public words and deeds is not omniscience about the heart. The \emph{Didache} 11 already keeps teaching and conduct as distinct tests, while Irenaeus III.3 identifies the public apostolic succession and preaching as the doctrinal rule.
\end{studybox}

\sectionguard
\subsection*{Romans 6: freedom becomes a gifted service\enspace\properrefs{Ep.}}

Romans 6:19--23 concludes an argument begun with baptism into Christ's death (6:1--11) and continued through the command not to let sin reign in the mortal body (6:12--18). Paul's slavery comparison is deliberately accommodated to ``the infirmity of your flesh.'' It is not an endorsement of human slavery or a complete definition of Christian freedom. Its purpose is to deny a supposed neutral zone in which the members serve no governing love.

\begin{threestable}{Pauline hinge}{Chrysostom's direct exegesis}{Augustinian doctrinal precision}
\latin{Exhibuistis} / \latin{exhibete} & Paul asks at least the same earnestness toward the immeasurably better master; the members are truly presented in embodied service. & Grace does not make exhortation fictitious: the healed will truly acts, although its new power and love are gifts.\\
\latin{Liberati a peccato, servi autem facti Deo} & Liberation, service to righteousness, and holiness are already possessed fruits rather than only distant promises. & Freedom is not independence from God; the works rewarded are God's workmanship enacted by the believer.\\
\latin{Fructum ... in sanctificationem, finem vero vitam aeternam} & Present sanctification gives evidence for the future end held in hope; shame and death mark the opposite service. & Eternal life is reward and grace because the good life crowned is itself grace (\emph{On Grace and Free Will} 18--21).\\
\latin{Stipendia ... mors; gratia ... vita aeterna} & Paul refuses parallel terms: death is sin's pay, but life is not described as wages earned from God. & Augustine makes the asymmetry decisive against both boasting and the denial that good works matter.\\
\end{threestable}

Anthony's \emph{Sermon for the Seventh Sunday after Pentecost}, sec.~13, follows the appointed Pauline sequence through liberation from sin, service to God, fruit unto sanctification, and eternal life. His preaching sharply retains Paul's unequal terms: death is the wage produced by sin, whereas eternal life is God's grace. Read beside Chrysostom and Augustine, this later saintly witness joins moral urgency to gratuity without making sanctification either passive or meritorious leverage over God.

The same Latin noun, \latin{fructus}, links the Epistle to the Gospel without making the two passages identical. In Matthew, fruit is evidence by which a teacher is known; in Romans, sanctification is the fruit of transferred allegiance and divine liberation. The juxtaposition yields a double conclusion: the fruit Christ requires is not a self-created credential, and the gift Paul proclaims is not fruitless permission.

\sectionguard
\subsection*{Collect and Gradual: governed, taught, and enlightened\enspace\properrefs{Coll., Grad.}}

No direct patristic commentary on the Collect was located. Its syntax first confesses that God's \latin{providentia} does not fail \latin{in sui dispositione}, then asks him to remove \latin{noxia} and grant \latin{profutura}. Thomas Aquinas's \emph{Summa theologiae} I, q.~22, a.~4 illuminates the distinction: unfailing providence includes free causes acting freely. The prayer is neither anxiety that God may lose control nor permission to ignore Paul's imperative and Christ's test.

The Gradual answers the question of whose teaching heals perception. Psalm 33:12 calls children to hear the fear of the Lord; verse 6 calls them to approach and be enlightened. Augustine's direct exposition is multilevel but exact. At section 15 he hears Christ as the teacher speaking through ministers. At section 9 he says the nations approach Christ by faith, longing, and the twofold charity of God and neighbor; he explicitly associates the approach with reception of the Crucified's Body and Blood. At section 17, the verses following the Gradual's first line move immediately from hearing to doing: keep the tongue from evil, depart from evil, and do good. True teaching is therefore neither disembodied information nor private suspicion; it draws the hearer toward Christ, charity, sacrament, and action.

\sectionguard
\subsection*{Offertory and Secret: deprivation, perfection, and common profit\enspace\properrefs{Off., Sec.}}

The Offertory is incompletely understood if Daniel 3:40 is isolated from Azariah's furnace prayer. In verses 38--39 there is no prince, prophet, leader, holocaust, sacrifice, oblation, incense, or place to offer first-fruits. The petitioner does not pretend those losses are irrelevant. Instead he asks that contrite soul and humble spirit be received, and then compares the desired acceptance to rams, bulls, and thousands of fat lambs. Verse 41 continues with wholehearted following; verse 42 grounds the plea in God's mercy. Material deprivation becomes neither contempt for liturgy nor confidence in interior sincerity by itself.

While expounding the living sacrifice of Romans 12, Chrysostom cites the furnace prayer and treats eyes, tongue, hands, feet, and hearing as members presented to God rather than retained for selfish use. Read beside Romans 6, this supplies related illumination for embodied self-offering.

Augustine's \emph{Letter 111}, secs.~3 and 8--9, supplies sustained direct ancient reception. Writing amid violence and captivity, he first invokes Azariah's confession to resist self-righteous complaint, then quotes the prayer's precise deprivation---no altar, priest, or material oblation---and its appeal that contrite heart and humble spirit be accepted as sacrifice. He applies the whole movement to captive Christians without treating their suffering as guilt or the absent liturgy as unimportant. Cyprian's earlier \emph{Letter 80}, secs.~2--3, gives related rather than exact illumination: for imprisoned confessors, it coordinates the broken and humble heart, Wisdom's furnace-and-offering imagery, and the three youths' fidelity. Together these witnesses keep contrition, bodily danger, material worship, and trust from collapsing into one another.

Hebrews 9:23--10:18 supplies the canonical boundary for the Secret's singular: Christ's self-offering is once for all. St. John Paul II, \emph{Ecclesia de Eucharistia} 12--13, citing the Council of Trent, identifies the sacrifice of the Cross and the Eucharistic sacrifice as one sacrifice: the Mass makes Calvary sacramentally present without adding to or multiplying it, and Christ joins the Church's spiritual offering to his own. This official formulation governs the relation between the one sacrifice, the present gifts, and the offerers.

The Secret's grammar provides a compact theology:

\begin{fourstable}{Clause}{Grammatical force}{Doctrinal illumination}{Distinction}
\latin{legalium differentiam hostiarum} & The legal victims are genuinely diverse. & Augustine, \emph{City of God} X.20, calls ancient sacrifices various signs of the one true sacrifice. & Diversity of figures does not mean several saving sacrifices parallel to Christ.\\
\latin{unius sacrificii perfectione sanxisti} & God's completed act is confessed before the petitions: one sacrifice gives perfection and definitive order. & \emph{Ecclesia de Eucharistia} 12--13, following Trent, says the Mass makes Calvary present without adding to or multiplying it; Augustine names Christ priest and victim. & The Church's gifts and self-offering are joined to Christ's; they are not additional redemptive sacrifices or coordinate victims.\\
\latin{sicut munera Abel, sanctifica} & God is asked to sanctify the present gifts as Abel's gifts were accepted. & Irenaeus IV.18 joins Abel's acceptance to righteousness, single-mindedness, and communion with the neighbor. & Moral disposition does not cause the objective sacrifice, but hypocrisy cannot profit by outward correctness.\\
\latin{quod singuli obtulerunt ... cunctis proficiat} & A singular relative clause gathers what individuals offered and asks that it profit everyone unto salvation. & Augustine's redeemed city is offered as one sacrifice through the High Priest (X.6). & The text grounds a personal-to-common movement, not a theory that every private intention has identical effects.\\
\end{fourstable}

\sectionguard
\subsection*{Communion and Postcommunion: the delivered still need cure\enspace\properrefs{Comm., Postcomm.}}

Psalm 30 moves from trust through urgent rescue to committal of the spirit. Augustine's \emph{Exposition on Psalm 30}, 1--4, first assigns the voice to the humbled Mediator and then to the people redeemed by his blood. He interprets ``bend down your ear'' as hearing Christ in humiliation and ``make haste to deliver me'' within the completion of his body and the final separation from sinners. Whatever one makes of each prosopological detail, the direct exposition blocks a complacent reading: Communion places the Church within Christ's petitioning voice.

The Postcommunion's verbs state both release and direction. The operation is God's; its recipients are \latin{nos}; it loosens them \latin{a nostris perversitatibus} and leads them \latin{ad ea quae sunt recta}. Medicine does not merely cancel an external penalty: it addresses perverse turnings and conducts the will toward right objects and acts.

\begin{sourcecard}{St. Ignatius and St. Thomas --- doctrinal illumination}
Ignatius, \emph{Ephesians} 20, calls the one bread broken in ecclesial unity the ``medicine of immortality'' and antidote against death. Thomas, \emph{Summa theologiae} III, q.~79, a.~6, says the Eucharist strengthens spiritual life as food and medicine and guards against sin through union with Christ and the Passion; his reply to the first objection retains the continuing mutability of free will.
\end{sourcecard}

The disputed-author \emph{De sacramentis} IV.6.27--28 offers a close liturgical parallel, joining Abel's gifts, Eucharistic remission, and repeated need for medicine.

The Collect and Postcommunion share a demonstrable grammatical direction. The first asks removal of \latin{noxia} and the gift of \latin{profutura}; the last asks release from \latin{perversitates} and guidance to \latin{recta}. Both place amendment under divine operation while retaining the recipients as the persons redirected. No checked witness joins the final prayer to the Gospel's tree image; that connection remains an editorial proposal.

\sectionguard
\section*{The Propers: Notable and Quotable}

{\small
\begin{itemize}[itemsep=0.35em,topsep=0.2em,parsep=0pt]
\item \textbf{Gospel, Matthew~7:15 --- New London broadside against Isaac Wilkins (4 May 1775).} ``A Wolf in Sheep's Clothing! A Traitor!'' stands above a hundred-pound reward and an order to seize Wilkins after his alleged betrayal of American liberties. Christ's pastoral warning becomes Revolutionary political typography and a public manhunt label; the broadside proves the phrase's polemical force, not the truth of its accusation.

\item \textbf{Gospel, Matthew~7:20 --- Rust College motto.} First president A.~C. McDonald made ``By Their Fruits Ye Shall Know Them'' the test of an education meant for long usefulness to students, their race, and the Church. A warning for discerning false prophets becomes an institutional outcome measure: alumni and service are the ``fruit'' by which a historically Black college founded in 1866 asks to be known.

\item \textbf{Epistle, Romans~6:23 --- Paula Poundstone, interview in \emph{Eugene Weekly} (7 April 2022).} Poundstone recalled seeing a billboard declare that the wages of sin are death and answering, ``after they take taxes out, it's just a tired feeling, really.'' Paul's ultimate wage becomes an ordinary paycheck diminished by withholding: a compact tax-season joke that turns death into post-payroll fatigue without pretending to interpret the verse.
\end{itemize}
}

\sectionguard
\section*{The Propers: Interpretive Possibilities}

\noindent\emph{The following exploratory proposals were developed in an AI-assisted editorial process. They are attributed to none of the cited authorities and claim no historical compositional intent. Each states its anchors and limits.}

\Needspace{8\baselineskip}
\noindent\textbf{The body cannot outsource its confession.\enspace\properrefs{Int., Ep., All., Gosp.}}
The Introit and Alleluia summon hands and voice to acclaim the universal King; Romans commands that bodily members change masters and bear sanctifying fruit; the Gospel tests religious voices by fruit and obedience. The mechanism is a progressive embodiment of confession: sound reaches hands, hands become members under a master, and members disclose their tree by fruit. Read one unit at a time, acclamation, moral conversion, and doctrinal testing can remain separate; together they propose that a community's operative habits are part of what its praise says. The intellectual and pastoral fruit is an institutional examen: budgets, incentives, protection of the vulnerable, and ordinary practices must not contradict voiced orthodoxy. The controlling limit is that works do not earn the gift of eternal life, institutional metrics cannot see the heart, and Augustine's hand/voice metaphor does not make every bodily action an unambiguous confession.

\Needspace{8\baselineskip}
\noindent\textbf{The wolf and Abel name rival liturgies of the flock.\enspace\properrefs{Gosp., Sec.}}
The Gospel's wolf borrows a sheep's exterior to consume the flock; the Secret asks that gifts be sanctified like Abel's; Abel is shepherd, offerer from the flock, and later the victim of fraternal violence. The mechanism is an inversion of the relation between religious appearance, flock, and appetite: one enters the flock as disguise and takes life, while the other gives from the flock and becomes vulnerable. Element-by-element reading asks whom the community should detect and what it should offer; the conjunction asks the more dangerous question of what its worship is making it desire. The fruit is self-discerning vigilance that tests protectors and offerers for predatory appetite before it is hidden by correct clothing. The limit is that Cain is not identified with Christ's wolf, Abel is not reduced to passive victimhood, and the Secret itself does not announce this narrative contrast.

\Needspace{8\baselineskip}
\noindent\textbf{What profits us is not permitted to stop with us.\enspace\properrefs{Coll., Sec., Postcomm.}}
The Collect asks providence to remove harm and give \latin{omnia nobis profutura}; the Secret asks that what individuals offer \latin{cunctis proficiat}; the Postcommunion asks medicinal operation to redirect the recipients toward what is right. The mechanism is an outward-turning circuit of benefit: gift received, gift offered, common profit, recipients redirected. A separate reading can leave providence as private protection and offering as private generosity; together they suggest that what truly benefits the Church tends to become shareable and to change the beneficiary into an agent of the good. The pastoral fruit is a criterion for discernment of resources and charisms: their ``profit'' includes the common body's healing and the giver's conversion. The limit is lexical and theological: \latin{profutura} and \latin{proficiat} are not the same lemma, ``us'' is already corporate, and usefulness alone does not define the good.

\Needspace{8\baselineskip}
\noindent\textbf{Two fires ask whether exposure ends or converts.\enspace\properrefs{Gosp., Off., Comm.}}
The Gospel sends the barren tree into fire; the Offertory extracts an offering from Azariah's prayer inside a furnace; Communion urgently asks the Lord to bend near and rescue. The mechanism is not one symbolic fire but a contrast between exposure that issues in judgment and ordeal from within which humble confession and petition remain possible. Read separately, fire can become only punishment and rescue only escape; the conjunction asks whether exposure has disclosed settled refusal or opened the sufferer to truthful offering and dependence. Its spiritual fruit is a discipline for crisis: seek conversion within trial without calling the trial good or presuming its outcome. The strongest limit is that Matthew and Daniel narrate different fires, the Offertory omits the furnace itself, and no one can infer from survival, suffering, or religious speech that judgment has been averted.

\Needspace{8\baselineskip}
\noindent\textbf{Medicine treats the tree, not merely the fruit report.\enspace\properrefs{Gosp., Sec., Postcomm.}}
The Gospel uses fruit to disclose the tree; the Secret places distinct offerings within Christ's one sacrifice for common salvation; the Postcommunion asks medicinal action to release recipients from perversity and lead them aright. The mechanism moves from diagnosis through sacrificial participation to treatment of the tree rather than cosmetic fruit replacement. Separate reading can make discernment a screening task---exclude bad fruit---while the prayers ask whether worship can heal disordered fruit-bearers. Its fruit is accountable hope: name corruption while seeking conversion of loves, habits, and direction. Sacramental medicine is not automatic reform, cancels neither safeguarding nor judgment, and Christ's warning remains where repentance is refused.
